\pagestyle{empty}
\tight

\begin{tblr}{
  width=\textwidth,
  colspec={|Q[c,m,1.9cm]|X[l,m]|},
  hlines, vlines,
  hspan=minimal,
  row{1}={bg=headblue},
  rowsep=1pt, colsep=3pt
}
\SetCell{c,m}\hfil\logo\hfil & \textbf{POLITEKNIK NEGERI MALANG}\par
\textbf{JURUSAN TEKNOLOGI INFORMASI}\par
\textbf{PROGRAM STUDI: D4 TEKNIK INFORMATIKA} \\
\SetCell[c=2]{c} \textbf{RENCANA PEMBELAJARAN SEMESTER (RPS)} & \\
\end{tblr}

\begin{longtblr}{
  width=\textwidth,
  colspec={|Q[l,m,1.9cm]|X[0.85,l,m]|X[1.45,l,m]|X[0.75,l,m]|X[1.15,l,m]|},
  hlines, vlines,
  hspan=minimal,
  rowsep=1pt, colsep=2pt,
  row{1,3}={bg=headgray,font=\bfseries}
}
MATA KULIAH & KODE & BOBOT (SKS)/jam & SEMESTER & TGL. PENYUSUNAN \\
Penjaminan Mutu Perangkat Lunak & RTI255006 & 2 SKS/4 jam & 5 & 13 Agustus 2025 \\
OTORISASI & Dosen Pengembang RPS & Koordinator RMK & \SetCell[c=2]{l} Ka PRODI & \\
 & Dian Hanifudin Subhi, S.Kom., M.Kom. & & \SetCell[c=2]{l}{\vspace{8mm}\par Dr. Ely Setyo Astuti, S.T., M.T.} & \\

\SetCell[r=12]{l} Capaian Pembelajaran (CP) & \SetCell[c=4]{l,bg=headgray,font=\bfseries} Capaian Pembelajaran Lulusan Program Studi (CPL-Prodi) & & & \\
 & CPL06 & \SetCell[c=3]{l} Mampu merancang, mengimplementasikan, menguji, melakukan deployment, memelihara, serta menjamin mutu perangkat lunak yang menjawab permasalahan dan memenuhi kebutuhan pengguna. & & \\
 & \SetCell[c=4]{l,bg=headgray,font=\bfseries} Capaian Pembelajaran mata kuliah (CPL-MK) & & & \\
 & CPMK0608 & \SetCell[c=3]{l} Mahasiswa mampu merancang strategi pengujian, melakukan verifikasi dan validasi, serta mengevaluasi jaminan mutu perangkat lunak berdasarkan kebutuhan, standar mutu, dan konteks pengembangan perangkat lunak. & & \\
 & \SetCell[c=4]{l,bg=headgray,font=\bfseries} Kemampuan akhir tiap tahapan belajar (Sub-CPMK) & & & \\
 & SCPMK0608-04201 & \SetCell[c=3]{l} Mahasiswa mampu menganalisis prinsip software quality assurance, standar mutu, tipe dan level pengujian, serta konteks pengujian dalam SDLC. & & \\
 & SCPMK0608-04202 & \SetCell[c=3]{l} Mahasiswa mampu merancang dokumen test plan, test scenario, test case, test monitoring, test control, dan test report yang sesuai dengan kebutuhan perangkat lunak. & & \\
 & SCPMK0608-04203 & \SetCell[c=3]{l} Mahasiswa mampu menerapkan dan mengevaluasi functional testing menggunakan tools yang sesuai untuk memverifikasi kesesuaian perangkat lunak terhadap requirement. & & \\
 & SCPMK0608-04204 & \SetCell[c=3]{l} Mahasiswa mampu merancang, menerapkan, dan mengevaluasi non-functional testing serta menyusun rekomendasi peningkatan kualitas perangkat lunak berdasarkan quality attributes. & & \\

\SetCell[r=2]{l} Penilaian & \SetCell[c=4]{l} *Nilai CPMK didapat dari agregasi nilai SUB-CPMK yang dibebankan & & & \\
 & \SetCell[c=4]{l}{\tiny\begin{tblr}{width=\linewidth,colspec={|Q[c,m,0.1\linewidth]|Q[c,m,0.16\linewidth]|X[l,m]|X[l,m]|X[l,m]|X[l,m]|X[l,m]|Q[c,m,0.08\linewidth]|},hlines,vlines,hspan=minimal,rowsep=1pt,colsep=0.7pt,row{1,2}={bg=headgray,font=\bfseries}}
Kode CPMK & Kode SCPMK & \SetCell[c=5]{c} Bobot per Bentuk Penilaian & & & & & Total Bobot \\
 & & Tugas & Kuis & UTS & PBL & UAS & \\
\SetCell[r=1]{c} \SetCell[r=4]{c} CPMK0608 & SCPMK0608-04201 & 4\%: SQA, standar mutu, tipe dan level pengujian & 7\%: konsep SQA dan perencanaan awal & 5\%: strategi pengujian & & & 16\% \\
\SetCell[r=3]{c}  & SCPMK0608-04202 & 7\%: test plan, scenario, case, monitoring, report & 3\%: dokumentasi pengujian & 10\%: strategi dan dokumentasi pengujian & 5\%: QA strategy awal & & 25\% \\
 & SCPMK0608-04203 & 18\%: functional, E2E, API, unit/integration testing & & & 10\%: functional testing proyek & 5\%: evaluasi hasil uji & 33\% \\
 & SCPMK0608-04204 & 6\%: non-functional dan performance testing & & & 10\%: strategi QA komprehensif & 10\%: rekomendasi kualitas & 26\% \\
\SetCell[c=7]{r}\textbf{Total Per Penilaian} & & & & & & \textbf{100\%} \\
\end{tblr}} & & & \\

Diskripsi Singkat Mata Kuliah & \SetCell[c=4]{l} Mata kuliah Penjaminan Mutu Perangkat Lunak memberikan pengetahuan, keterampilan, dan pengalaman terarah untuk merancang strategi penjaminan mutu perangkat lunak melalui pengujian, verifikasi, validasi, dokumentasi pengujian, functional testing, non-functional testing, serta penyusunan rekomendasi peningkatan kualitas perangkat lunak sesuai standar industri dan kebutuhan pengguna. & & & \\
Bahan Kajian / Materi Pembelajaran & \SetCell[c=4]{l}\begin{enumerate}\item Pengantar SQA dan standar ISO.\item Pendekatan, tipe, dan level pengujian.\item Test planning, test scenario, test case, dan traceability.\item Test monitoring, test control, test completion, dan test reporting.\item Functional testing.\item Non-functional testing.\item Strategi QA komprehensif.\end{enumerate} & & & \\
\SetCell[r=2]{l} Pustaka & \SetCell[c=4]{l} \textbf{Utama:}\begin{enumerate}\item Sommerville, I. 2016. \textit{Software Engineering}, 10th Edition. Pearson.\item Myers, G. J., Sandler, C., and Badgett, T. 2011. \textit{The Art of Software Testing}, 3rd Edition. Wiley.\item Pressman, R. S. and Maxim, B. R. 2019. \textit{Software Engineering: A Practitioner's Approach}, 9th Edition. McGraw-Hill.\end{enumerate} & & & \\
 & \SetCell[c=4]{l} \textbf{Pendukung:}\begin{enumerate}\item ISO/IEC/IEEE 29119 Software Testing Standard.\item ISO/IEC 25010 Systems and Software Quality Requirements and Evaluation.\item ISTQB Certified Tester Foundation Level Syllabus.\end{enumerate} & & & \\
Media Pembelajaran & \SetCell[c=2]{l} \textbf{Software:} Office, IDE/text editor, Git, Postman/Newman, Selenium/Cypress, JMeter/k6, unit testing framework. & & \SetCell[c=2]{l} \textbf{Hardware:} Komputer/laptop, proyektor. & \\
Nama Dosen Pengampu & \SetCell[c=4]{l}\begin{enumerate}\item Dian Hanifudin Subhi, S.Kom., M.Kom.\item Dika Rizky Yunianto, S.Kom., M.Kom.\item M. Hasyim Ratsanjani, S.Kom., M.Kom.\end{enumerate} & & & \\
Matakuliah Syarat & \SetCell[c=4]{l} - & & & \\
\end{longtblr}

\begin{longtblr}{
  width=\textwidth,
  colspec={|Q[c,m,0.06\textwidth]|X[1.35,l,m]|X[1.08,l,m]|X[1.16,l,m]|X[0.7,l,m]|X[1.24,l,m]|X[1.06,l,m]|X[0.98,l,m]|Q[c,m,0.05\textwidth]|},
  hlines, vlines,
  hspan=minimal,
  rowhead=2,
  row{1,2}={bg=headgreen,font=\bfseries},
  rowsep=1pt, colsep=1.5pt
}
Mgg.\par Ke & Kemampuan Akhir Yang Direncanakan (Sub-CP-MK) & Bahan kajian (Materi Pembelajaran) & Bentuk dan Metode Pembelajaran & Estimasi Waktu & Pengalaman Belajar Mahasiswa & Kriteria \& Bentuk Penilaian & Indikator Penilaian & Bobot Penilaian (\%) \\
(1) & (2) & (3) & (4) & (5) & (6) & (7) & (8) & (9) \\
1 & \SetCell[r=2]{l} SCPMK0608-04201: Mahasiswa mampu menganalisis prinsip software quality assurance, standar mutu, tipe dan level pengujian, serta konteks pengujian dalam SDLC. & Pengantar SQA; software quality; ISO/IEC 25010; peran QA dalam SDLC. & Modalitas: Blended Learning. Bentuk: Luring/Daring. Metode: CTL, studi kasus, diskusi kelompok. Penugasan: Tugas 1 analisis standar ISO dan QA. & 1 x 3 x 50' tatap muka dan tugas terstruktur; 1 x 1 x 50' tugas kelompok. & Mahasiswa menjelaskan pentingnya QA, mengidentifikasi standar mutu, dan menganalisis masalah mutu pada kasus perangkat lunak. & Kriteria: rubrik kriteria penilaian. Bentuk: diskusi dan Tugas 1. Mengacu Rubrik R1. & Ketepatan analisis masalah mutu; kesesuaian standar; kejelasan argumentasi. & 2 \\
2 & & Pendekatan pengujian; manual vs automated; black-box, white-box, gray-box; unit, integration, system, acceptance testing. & Modalitas: Blended Learning. Metode: CTL dan Problem Based Learning. Penugasan: Tugas 2 identifikasi tipe dan level pengujian. & 1 x 2 x 50' tatap muka; 1 x 2 x 50' tugas kelompok. & Mahasiswa membandingkan pendekatan pengujian dan menentukan tipe/level pengujian untuk requirement tertentu. & Kriteria: ketepatan dan penguasaan. Bentuk: diskusi dan Tugas 2. Mengacu Rubrik R2. & Ketepatan klasifikasi; alasan pemilihan; keterkaitan dengan SDLC. & 2 \\
3 & SCPMK0608-04202: Mahasiswa mampu merancang dokumen test plan, test scenario, dan test case. & Test plan; test scenario; test case; requirement traceability. & Modalitas: Blended Learning. Metode: Project Based Learning dan studi kasus. Penugasan: Tugas 3 penyusunan test plan, scenario, dan case. & 1 x 2 x 50' tatap muka; 1 x 2 x 50' tugas kelompok. & Mahasiswa menyusun artefak perencanaan pengujian untuk aplikasi sederhana. & Kriteria: rubrik kriteria penilaian. Bentuk: Tugas 3. Mengacu Rubrik R4. & Kelengkapan test plan; logika scenario; detail test case; traceability. & 4 \\
4 & SCPMK0608-04201 dan SCPMK0608-04202 & Materi minggu 1--3 dan artefak awal proyek kelompok. & Evaluasi: Kuis 1 berbasis checkpoint PBL kelompok. & 1 x 2 x 50' evaluasi dan review proyek; 1 x 2 x 50' penyempurnaan artefak kelompok. & Mahasiswa mempresentasikan konteks QA proyek, pemetaan ISO/quality model, ruang lingkup pengujian, risiko awal, dan traceability awal. & Bentuk penilaian: Kuis 1. Mengacu Rubrik R3. & Ketepatan konteks QA; relevansi standar; kejelasan testing scope; prioritas risiko; traceability awal. & 10 \\
5 & \SetCell[r=2]{l} SCPMK0608-04202: Mahasiswa mampu merancang dokumen test plan, test scenario, test case, test monitoring, test control, dan test report yang sesuai dengan kebutuhan perangkat lunak. & Review test plan; review scenario; review test case; best practices dokumentasi pengujian. & Modalitas: Blended Learning. Metode: Peer Review dan Project Based Learning. Penugasan: Tugas 4 review dan presentasi artefak. & 1 x 2 x 50' tatap muka; 1 x 2 x 50' tugas kelompok. & Mahasiswa melakukan peer review, mempresentasikan temuan, dan memperbaiki dokumen pengujian. & Bentuk penilaian: Tugas 4 dan presentasi. Mengacu Rubrik R5. & Kedalaman review; rekomendasi; kualitas presentasi; revisi artefak. & 2 \\
6 & & Test monitoring; test control; test completion; defect summary; test report. & Modalitas: Blended Learning. Metode: CTL dan Project Based Learning. Penugasan: Tugas 5 simulasi monitoring dan reporting. & 1 x 2 x 50' tatap muka; 1 x 2 x 50' tugas kelompok. & Mahasiswa memonitor progres pengujian, mengendalikan risiko, dan menyusun laporan pengujian. & Bentuk penilaian: Tugas 5. Mengacu Rubrik R6. & Ketepatan metrik; kelengkapan laporan; kontrol risiko; rekomendasi tindak lanjut. & 2 \\
7 & SCPMK0608-04203: Mahasiswa mampu menerapkan dan mengevaluasi functional testing. & Functional testing; test design technique; UAT planning; user story testing. & Modalitas: Blended Learning. Metode: Hands-on Practice dan Role Playing. Penugasan: Tugas 6 functional testing dan Tugas 7 UAT. & 1 x 2 x 50' tatap muka; 1 x 2 x 50' tugas kelompok. & Mahasiswa merancang pengujian fungsional dan melakukan simulasi UAT sebagai user dan tester. & Bentuk penilaian: Tugas 6 dan Tugas 7. Mengacu Rubrik R8 dan R9. & Coverage requirement; validitas test data; kualitas UAT; bukti hasil uji. & 6 \\
8 & UTS untuk SCPMK0608-04201 dan SCPMK0608-04202 & Materi minggu 1--7 dan dokumen milestone proyek kelompok. & Evaluasi: UTS berbasis milestone review PBL kelompok. & 1 x 3 x 50' evaluasi praktik dan defense proyek. & Mahasiswa mempresentasikan dan mempertahankan test plan, test scenario, test case, traceability matrix, serta monitoring/control plan proyek. & Bentuk penilaian: UTS. Mengacu Rubrik R7. & Kelengkapan dokumen; konsistensi strategi; traceability; monitoring/control plan; kualitas defense kelompok. & 15 \\
9 & \SetCell[r=4]{l} SCPMK0608-04203: Mahasiswa mampu menerapkan dan mengevaluasi functional testing menggunakan tools yang sesuai untuk memverifikasi kesesuaian perangkat lunak terhadap requirement. & End-to-end testing; test flow; Selenium/Cypress; test script. & Modalitas: Blended Learning. Metode: Hands-on Practice dengan tools. Penugasan: Tugas 8 implementasi E2E testing. & 1 x 2 x 50' tatap muka; 1 x 2 x 50' tugas kelompok. & Mahasiswa membuat skenario E2E, menulis test script, dan mendokumentasikan hasil. & Bentuk penilaian: Tugas 8. Mengacu Rubrik R10. & Kebenaran flow E2E; kualitas script; hasil eksekusi; dokumentasi evidence. & 3 \\
10 & & API testing; Postman; assertion; Newman; test suite. & Modalitas: Blended Learning. Metode: Hands-on Practice dengan Postman/Newman. Penugasan: Tugas 9 automated API testing suite. & 1 x 2 x 50' tatap muka; 1 x 2 x 50' tugas kelompok. & Mahasiswa merancang request, assertion, collection, environment, dan automated run. & Bentuk penilaian: Tugas 9. Mengacu Rubrik R11. & Ketepatan request; kelengkapan assertion; struktur collection; bukti eksekusi otomatis. & 3 \\
11 & & Unit testing; integration testing; system testing; TDD approach. & Modalitas: Blended Learning. Metode: Coding Practice dan TDD. Penugasan: Tugas 10 unit dan integration test. & 1 x 2 x 50' tatap muka; 1 x 2 x 50' tugas kelompok. & Mahasiswa menulis test code, menjalankan test suite, dan mengevaluasi hasil. & Bentuk penilaian: Tugas 10. Mengacu Rubrik R12. & Kualitas test code; coverage fungsi kritis; penerapan TDD; integrasi kode. & 3 \\
12 & & Review unit/integration test; code review; coverage analysis; test maintainability. & Modalitas: Blended Learning. Metode: Code Review dan Peer Assessment. Penugasan: Tugas 11 evaluasi test code. & 1 x 2 x 50' tatap muka; 1 x 2 x 50' tugas kelompok. & Mahasiswa melakukan review test code, menganalisis coverage, dan memberi rekomendasi perbaikan. & Bentuk penilaian: Tugas 11. Mengacu Rubrik R13. & Ketepatan review; pemahaman coverage; rekomendasi refactoring test; kualitas feedback. & 3 \\
13 & \SetCell[r=2]{l} SCPMK0608-04204: Mahasiswa mampu merancang, menerapkan, dan mengevaluasi non-functional testing serta menyusun rekomendasi peningkatan kualitas perangkat lunak berdasarkan quality attributes. & Non-functional testing; performance; security; usability; maintainability; ISO/IEC 25010 quality attributes. & Modalitas: Blended Learning. Metode: CTL dan Case Study Analysis. Penugasan: Tugas 12 analisis kebutuhan non-functional testing. & 1 x 2 x 50' tatap muka; 1 x 2 x 50' tugas kelompok. & Mahasiswa mengidentifikasi quality attributes dan merancang strategi non-functional testing. & Bentuk penilaian: Tugas 12. Mengacu Rubrik R14. & Identifikasi atribut kualitas; relevansi strategi; prioritas risiko; keterkaitan kebutuhan pengguna. & 3 \\
14 & & Load testing; performance testing; stress testing; volume testing; JMeter/k6; performance metrics. & Modalitas: Blended Learning. Metode: Hands-on Practice dengan JMeter/k6. Penugasan: Tugas 13 performance testing scenario. & 1 x 2 x 50' tatap muka; 1 x 2 x 50' tugas kelompok. & Mahasiswa membuat skenario performa, menjalankan pengujian, membaca metrik, dan menyusun rekomendasi. & Bentuk penilaian: Tugas 13. Mengacu Rubrik R15. & Kesesuaian skenario beban; validitas metrik; analisis bottleneck; rekomendasi. & 3 \\
15--16 & SCPMK0608-04201--04204: Integrasi semua kemampuan dalam proyek komprehensif QA. & Proyek akhir; QA strategy; test planning; functional testing; non-functional testing; reporting; presentation. & Modalitas: Blended Learning. Metode: Project Based Learning dan mentoring. Penugasan: PBL proyek akhir strategi QA lengkap. & 2 x 2 x 50' tatap muka; 2 x 2 x 50' tugas kelompok. & Mahasiswa menyusun strategi QA, mengimplementasikan pengujian, membuat laporan, dan mempresentasikan hasil. & Bentuk penilaian: PBL/proyek akhir. Mengacu Rubrik R16. & Integrasi strategi; implementasi; evidence pengujian; dokumentasi; presentasi; rekomendasi mutu. & 25 \\
17 & UAS untuk SCPMK0608-04203 dan SCPMK0608-04204 & Materi minggu 1--16 dan hasil akhir proyek kelompok. & Evaluasi: UAS berbasis validasi akhir PBL kelompok. & 1 x 3 x 50' evaluasi praktik dan presentasi akhir. & Mahasiswa mendemonstrasikan implementasi QA lengkap, evidence functional dan non-functional testing, analisis defect, serta rekomendasi peningkatan mutu. & Bentuk penilaian: UAS. Mengacu Rubrik R17. & Kelengkapan evidence; ketepatan analisis defect; kualitas rekomendasi; konsistensi dokumentasi; kualitas presentasi akhir. & 15 \\
\end{longtblr}

